\documentclass[11pt]{amsart}
\usepackage{amssymb,amsmath,booktabs}
\usepackage[margin=1.1in]{geometry}
\usepackage[hidelinks]{hyperref}

\newtheorem{theorem}{Theorem}
\newtheorem{corollary}[theorem]{Corollary}
\newtheorem{lemma}[theorem]{Lemma}
\theoremstyle{remark}
\newtheorem{remark}[theorem]{Remark}

\newcommand{\Z}{\mathbb{Z}}

\title{The determinant of MATLAB's odd-order magic squares}
\author{Brayden Sanders}
\address{Hot Springs, Arkansas, USA}
\email{}
\date{October 30, 2026}

\begin{document}

\begin{abstract}
For odd $n$ we give a closed form for the determinant of the $n\times n$ magic square returned by MATLAB's
\texttt{magic(n)} (the de la Loub\`ere, or siamese, construction):
\[
\det \mathrm{magic}(n) \;=\; (-1)^{\frac{n-1}{2}}\, n^{\,n-1}\,\frac{n^2+1}{2}\prod_{O}\bigl(n^{|O|}-(-1)^{|O|}\bigr),
\]
where $O$ runs over the orbits of $k\mapsto 2k$ on $\Z_n\setminus\{0\}$. In particular the determinant never
vanishes, and $\det\mathrm{magic}(n)/n^{n-1}$ is an integer coprime to $n$. This proves the divisibility by $n^{n-1}$
observed in OEIS A160523. The proof subtracts the mean entry and works in the discrete Fourier basis, where the
centered matrix becomes a diagonal matrix plus a weighted permutation. The same argument shows that every odd-order
square built from two linear Latin squares with unit coefficients is nonsingular.
\end{abstract}

\maketitle

\section{Introduction}

MATLAB's \texttt{magic(n)} returns an $n\times n$ magic square. For odd $n$ it uses the de la Loub\`ere (siamese)
construction. In matrix form, with rows $I$ and columns $J$ running over $1,\dots,n$,
\[
\mathrm{magic}(n)_{IJ} \;=\; n\cdot\bigl((I+J-\tfrac{n+3}{2}) \bmod n\bigr)+\bigl((I+2J-2)\bmod n\bigr)+1 .
\]
Moler \cite{moler} discussed the linear algebra of these squares. Numerically, odd orders have full rank, orders
divisible by $4$ have rank $3$, and the remaining even orders have rank $n/2+2$. The OEIS entry A160523 \cite{oeis}
lists $a(n)=\det\mathrm{magic}(n)$. There Robert Israel noted (2015) that $a(n)$ is divisible by $n^{n-2}$ for odd
$n$ and appears to be divisible by $n^{n-1}$. Manjunath and Ambikasaran \cite{ma} approximated the eigenvalues of
the odd-order squares with proven error bounds; nonsingularity follows from those bounds, although it is not stated
there. We were unable to find an exact formula for $a(n)$ in the literature.

\begin{theorem}\label{thm:main}
Let $n$ be odd. Then
\[
\det \mathrm{magic}(n) = (-1)^{\frac{n-1}{2}}\, n^{\,n-1}\,\frac{n^2+1}{2}\prod_{O}\bigl(n^{|O|}-(-1)^{|O|}\bigr),
\]
where the product runs over the orbits $O$ of the doubling map $k\mapsto 2k$ on $\Z_n\setminus\{0\}$. For $n=1$ the
product is empty.
\end{theorem}

\begin{corollary}
For odd $n$, $\mathrm{magic}(n)$ is nonsingular. Moreover $\det\mathrm{magic}(n)/n^{n-1}$ is an integer coprime to
$n$, so for each prime $p\mid n$ the exact power of $p$ dividing the determinant is $p^{(n-1)v_p(n)}$.
\end{corollary}

\begin{proof}
Each factor $n^{L}-(-1)^{L}$ is nonzero and congruent to $\pm1 \pmod n$, and $\gcd\bigl(\tfrac{n^2+1}{2},n\bigr)=1$.
\end{proof}

For example, for $n=7$ the orbits are $\{1,2,4\}$ and $\{3,6,5\}$, so
$\det\mathrm{magic}(7) = -7^6\cdot 25\cdot (7^3+1)^2 = -348\,052\,801\,600$.

\section{Proof of Theorem~\ref{thm:main}}

Index rows and columns by $i=I-1$ and $j=J-1$ in $\Z_n$, and put $h=\frac{n-1}{2}$. Then
\[
M_{ij} = n\bigl((i+j-h)\bmod n\bigr) + \bigl((i+2j+1)\bmod n\bigr) + 1 .
\]
Let $f(x)=(x\bmod n)-h$. This is an $n$-periodic function with $\sum_{x\in\Z_n}f(x)=0$.

\begin{lemma}[Measuring from the center]
$M = N + \frac{n^2+1}{2}\,\mathbf{1}\mathbf{1}^{\mathsf T}$, where $N_{ij}=n\,f(i+j-h)+f(i+2j+1)$.
\end{lemma}

\begin{proof}
Write $(x\bmod n)=f(x)+h$. The constant term is then $nh+h+1=(n+1)\frac{n-1}{2}+1=\frac{n^2+1}{2}$.
\end{proof}

Let $\omega=e^{2\pi i/n}$, let $e_k(j)=\omega^{kj}$ for $k\in\Z_n$, and let
$\hat f(k)=\sum_{x\in\Z_n}f(x)\,\omega^{kx}$. For $k\neq0$,
\begin{equation}\label{eq:saw}
\hat f(k)=\sum_{x=0}^{n-1}x\,\omega^{kx}=\frac{n}{\omega^k-1},
\end{equation}
because $(z-1)\sum_{x=0}^{n-1}xz^x=n$ whenever $z^n=1\neq z$.

\begin{lemma}[Fourier action]
For every $k\in\Z_n$,
\[
N e_k = a_k\, e_{-k} + b_{k/2}\, e_{-k/2},\qquad a_k=n\,\omega^{kh}\hat f(k),\quad b_k=\omega^{-k}\hat f(k),
\]
where $k/2$ means $k\cdot 2^{-1}$ in $\Z_n$ (recall $n$ is odd).
\end{lemma}

\begin{proof}
Substitute $x=i+j-h$ in $\sum_j f(i+j-h)\,\omega^{kj}$ to get $\omega^{kh}\hat f(k)\,\omega^{-ki}$. Substitute
$y=i+2j+1$, so that $j=(y-i-1)/2$, in $\sum_j f(i+2j+1)\,\omega^{kj}$ to get $\hat f(k/2)\,\omega^{-k/2}\,\omega^{-(k/2)i}$.
\end{proof}

Since $\hat f(0)=0$, the lemma gives $N e_0=0$, and so $M e_0=s\,e_0$ with $s=n\frac{n^2+1}{2}$. For $k\neq0$ we have
$\mathbf{1}^{\mathsf T}e_k=0$, so $M e_k=N e_k$. Both $\mathbb{C}e_0$ and $V=\mathrm{span}\{e_k:k\neq0\}$ are therefore
invariant, and
\[
\det M = s\cdot \det\bigl(N|_V\bigr).
\]

On $V$, define $R e_k=e_{-k}$, $D e_k=a_k e_k$ and $U e_k=b_{k/2}\,e_{k/2}$. The lemma says exactly that
$N|_V=R\,(D+U)$. The involution $k\mapsto -k$ has no fixed points on $\Z_n\setminus\{0\}$, so
$\det R=(-1)^{(n-1)/2}$.

The halving map $k\mapsto k/2$ permutes $\Z_n\setminus\{0\}$, and its cycles are the orbits $O$ of doubling.
Hence $D+U$ is block diagonal over the orbits. Each block is a diagonal matrix plus one weighted cyclic permutation
of length $L=|O|\geq2$. In such a block the only permutations in the support are the identity and the full cycle,
so its determinant is
\[
\prod_{k\in O}a_k-(-1)^{L}\prod_{k\in O}b_k
= \Bigl(\prod_{k\in O}\hat f(k)\Bigr)\Bigl(n^{L}\omega^{hS_O}-(-1)^{L}\omega^{-S_O}\Bigr),
\qquad S_O=\sum_{k\in O}k .
\]
If $O=\{k,2k,\dots,2^{L-1}k\}$ then $2^{L}k\equiv k$, so $S_O=(2^{L}-1)k\equiv0\pmod n$. Both phases are therefore
$1$, and the block determinant is $(n^{L}-(-1)^{L})\prod_{k\in O}\hat f(k)$.

Finally, by \eqref{eq:saw},
\[
\prod_{k\neq0}\hat f(k)=\frac{n^{n-1}}{\prod_{k=1}^{n-1}(\omega^k-1)}=\frac{n^{n-1}}{n}=n^{n-2},
\]
using $\prod_{k=1}^{n-1}(1-\omega^k)=n$ and the fact that $n-1$ is even. Collecting the factors,
\[
\det M = n\,\tfrac{n^2+1}{2}\cdot(-1)^{\frac{n-1}{2}}\cdot n^{n-2}\prod_O\bigl(n^{|O|}-(-1)^{|O|}\bigr),
\]
which is the theorem. \qed

\section{A family of nonsingular squares}

The proof uses only the fact that each of the two digit layers is linear in $(i,j)$ with unit coefficients.

\begin{theorem}\label{thm:family}
Let $n\geq3$ be odd, let $a_1,b_1,a_2,b_2$ be units of $\Z_n$, and let $c_1,c_2\in\Z_n$. With $\rho(x)=x\bmod n$,
the matrix
\[
M_{ij}=n\,\rho(a_1i+b_1j+c_1)+\rho(a_2i+b_2j+c_2)+1
\]
is nonsingular.
\end{theorem}

\begin{proof}
All row sums are equal because $b_1,b_2$ are units, so $M e_0=s e_0$ with $s\neq0$ and $V$ is invariant as before.
Substituting as in the Fourier lemma gives $N e_k=A_k e_{\sigma(k)}+B_k e_{\tau(k)}$, where $\sigma(k)=-a_1b_1^{-1}k$
and $\tau(k)=-a_2b_2^{-1}k$. By \eqref{eq:saw},
\[
|A_k|=\frac{n^2}{|\omega^{k/b_1}-1|}\geq\frac{n^2}{2},\qquad
|B_k|=\frac{n}{|\omega^{k/b_2}-1|}\leq\frac{n}{2\sin(\pi/n)} .
\]
Write $N|_V=P_\sigma(D_A+W)$, where $W e_k=B_k e_{\pi(k)}$ and $\pi=\sigma^{-1}\tau$. Each cycle of $\pi$ of length
$L$ contributes $\prod A_k-(-1)^L\prod B_k$, and
\[
\Bigl|\prod A_k\Bigr| \big/ \Bigl|\prod B_k\Bigr| \;\geq\; \bigl(n\sin(\pi/n)\bigr)^{L} > 1
\]
because $n\sin(\pi/n)\geq 3\sin(\pi/3)>2$ for $n\geq3$. If $\pi$ fixes $k$, the factor $A_k+B_k$ is nonzero for the same
reason. So no factor vanishes.
\end{proof}

Theorem~\ref{thm:main} is the case $(a_1,b_1,c_1,a_2,b_2,c_2)=(1,1,-h,1,2,1)$, where the two permutations are
$k\mapsto-k$ and $k\mapsto-k/2$ and the phases cancel orbit by orbit.

\section{Verification}

We computed $\det\mathrm{magic}(n)$ exactly, by fraction-free Gaussian elimination over the integers, for every odd
$n\le101$. Every value agrees with Theorem~\ref{thm:main}, and the formula also reproduces the ten odd-index terms
listed in A160523 ($n\le19$). A separate implementation, written independently from MATLAB's \texttt{magic.m} source,
confirmed the same range. Theorem~\ref{thm:family} was tested on $77{,}904$ squares: all parameter choices for
$n=5$ and $7$, and $2{,}000$ random choices each for $n=9,15,21,25$. No zero determinant occurred.

\begin{table}[h]
\centering
\begin{tabular}{rll}
\toprule
$n$ & orbit lengths of $k\mapsto2k$ & $\det\mathrm{magic}(n)/n^{n-1}$ \\
\midrule
3 & 2 & $-40$ \\
5 & 4 & $8112$ \\
7 & 3, 3 & $-2958400$ \\
9 & 6, 2 & $1743123200$ \\
11 & 10 & $-1582182900600$ \\
13 & 12 & $1980337235410800$ \\
15 & 4, 4, 2, 4 & $-3283944692869234688$ \\
\bottomrule
\end{tabular}
\end{table}

\begin{remark}
The formula depends on the arithmetic of $n$ only through the orbit lengths of doubling, that is, through the
multiplicative orders of $2$ modulo the divisors of $n$. When $2$ is a primitive root mod a prime $n$, there is a
single orbit and $\det\mathrm{magic}(n)=(-1)^{(n-1)/2}n^{n-1}\frac{n^2+1}{2}\,(n^{n-1}-1)$.
\end{remark}

\subsection*{Acknowledgment}
The derivation, the computations and the drafting of this note were carried out with the assistance of an AI system
(Claude, Anthropic). The author checked the argument and the computations and takes responsibility for the content.

\begin{thebibliography}{9}
\bibitem{moler} C. Moler, \emph{Magic Squares, Part 3, Linear Algebra}, Cleve's Corner (MathWorks blog),
November 12, 2012. \url{https://blogs.mathworks.com/cleve/2012/11/12/magic-squares-part-3-linear-algebra/}
\bibitem{oeis} OEIS Foundation, \emph{The On-Line Encyclopedia of Integer Sequences}, sequence A160523.
\url{https://oeis.org/A160523}
\bibitem{ma} H. Manjunath and S. Ambikasaran, \emph{Spectrum of MATLAB's magic squares},
arXiv:2205.15563 (2022).
\bibitem{staab} P. Staab et al., \emph{The magic of permutation matrices: categorizing, counting and eigenspectra
of magic squares}, arXiv:1007.3239 (2010).
\end{thebibliography}

\end{document}
